\documentclass[10pt,journal,compsoc,twocolumn]{article}

\usepackage[margin=1.8cm,top=2.0cm,bottom=2.0cm]{geometry}
\usepackage{amsmath,amssymb,amsfonts}
\usepackage{algorithmic}
\usepackage{graphicx}
\usepackage{textcomp}
\usepackage{xcolor}
\usepackage{booktabs}
\usepackage{listings}
\usepackage{microtype}
\usepackage{hyperref}
\usepackage{enumitem}
\usepackage{cite}
\usepackage{titlesec}
\usepackage{fancyhdr}
\usepackage{CJKutf8}

\newcommand{\ja}[1]{\begin{CJK*}{UTF8}{min}#1\end{CJK*}}

\hypersetup{
    colorlinks=true,
    linkcolor=blue!70!black,
    citecolor=blue!70!black,
    urlcolor=blue!70!black
}

\lstset{
    basicstyle=\ttfamily\scriptsize,
    breaklines=true,
    frame=single,
    backgroundcolor=\color{gray!6},
    keywordstyle=\color{blue!80!black}\bfseries,
    commentstyle=\color{green!50!black}\itshape,
    stringstyle=\color{red!70!black},
    numbers=none,
    showstringspaces=false,
    tabsize=2
}

\title{\LARGE \textbf{TAD-Native Declarative Applications in BTRON:\\The NITRO/FORM Framework for Structured Document-Centric Interaction in SPEC 3.20}}

\author{
    \textbf{Maxim Sokhatsky (Namdak Tonpa)}\\[0.3em]
    \small \textit{Synrc Research Center / Groupoid Infinity Cleanroom Working Group}\\
    \small \textit{Kyiv, Ukraine}\\
    \small \texttt{maxim@synrc.com}
}
\date{August 2026}

\pagestyle{fancy}
\fancyhf{}
\rhead{\small \textit{TAD-Native Declarative Applications in BTRON}}
\lhead{\small \textit{Synrc Research Report}}
\cfoot{\thepage}

\begin{document}

\maketitle

\begin{abstract}
The historical divergence between document interchange formats and interactive software applications has led to pervasive complexity in modern graphical user interface (GUI) stacks. In contrast, Ken Sakamura's \textbf{BTRON} architecture unifies documents, applications, and operating system data structures into a single polymorphic medium: the \textbf{TRON Application Databus (TAD)}. In this paper, we present \textbf{SYNRC NITRO/FORM}, a lightweight, declarative forms and enterprise application framework designed as a strict, non-invasive extension of the \textbf{BTRON3 SPEC 3.20} standard. NITRO/FORM models business forms as first-class TAD Real Bodys (\ja{実身}), encapsulating dynamic UI controls, reactive data binding, and validation schemas within standard application-range \emph{fusen} (\ja{付箋}) segments. We formalize the tripartite form algebra (Header, Body, Footer), define a bidirectional transcompilation bridge between TAD forms and modern semantic HTML5, and prove the Graceful Degradation Invariant guaranteeing backward compatibility with legacy BTRON renderers. We demonstrate experimental performance across POSIX, QEMU VirtIO, and baremetal ARM Cortex-A architectures, showing that document-native reactive applications achieve high determinism and sub-millisecond layout latencies.
\end{abstract}

\noindent\textbf{\small Keywords}---BTRON, TRON Application Databus (TAD), NITRO/FORM, Declarative UI, Real Bodys, Fusen, Parts Manager, X-Forms, Cleanroom Architecture.

%-------------------------------------------------------------------------
\section{Introduction and Theoretical Motivation}
\label{sec:intro}

In contemporary software engineering, a deep architectural dichotomy exists between static document markup (e.g., PDF, static HTML) and imperative application runtimes (e.g., JavaScript DOM manipulation frameworks, desktop widget toolkits). This division forces distributed systems to maintain complex object-relational mappings, fragile serialization layers, and heavyweight client-side runtime engines.

When Ken Sakamura formulated the \textbf{TRON} architecture in 1984, he proposed a fundamentally distinct model for personal and workstation computing: the **TRON Application Databus (TAD)**. Rather than treating applications as opaque binary executables that read and write external files, BTRON treats every entity in the system as a structured, hyper-linked TAD document residing in a persistent Real Body / Virtual Body (\ja{実身}/\ja{仮身}) graph.

Under BTRON3 SPEC 3.20, TAD establishes a universal container format capable of embedding multilingual text, 2D vector geometry, raster bitmapped images, sound, and executable Virtual Bodys into a coherent, stream-oriented record sequence.

However, classic BTRON specifications did not formalize a high-level, declarative framework for transactional business forms, structured data entry, and reactive schema validation (analogous to modern X-Forms or declarative component models).

This paper presents the design and formalization of **SYNRC NITRO/FORM**, a lightweight, document-centric application framework engineered for BTRON3 SPEC 3.20. NITRO/FORM achieves three primary technical goals:
\begin{enumerate}[leftmargin=*]
    \item \textbf{Strict Legal Extensibility:} Zero modifications to core $\mu$ITRON system calls, kernel structures, or standard TAD segment IDs, utilizing application-range \emph{fusen} tags.
    \item \textbf{Declarative Triple Structure:} Formalizing business forms into an algebra of Header, Body, and Footer regions bound to underlying data models.
    \item \textbf{Bidirectional Interchange:} Establishing a lossless bridge between native TAD forms and semantic HTML5 documents via the companion TAD Browser subsystem.
\end{enumerate}

%-------------------------------------------------------------------------
\section{TAD Segment Architecture and Fusen Mechanics}
\label{sec:tad_foundations}

\subsection{TAD Record Organization}
A standard TAD document consists of a stream of variable-length primitive segments, each identified by a 16-bit type tag ($T_{\text{tag}}$) and a length descriptor ($L_{\text{len}}$):

\begin{equation}
\text{Segment} = \langle T_{\text{tag}}, L_{\text{len}}, \mathcal{D}_{\text{payload}} \rangle
\end{equation}

Standard system segments defined in SPEC 3.20 include:
\begin{itemize}[leftmargin=*]
    \item \texttt{TS\_INFO}: Document metadata, author, and timestamp.
    \item \texttt{TS\_TEXT}: Multi-lingual TRON Code text stream.
    \item \texttt{TS\_FIG}: Vector graphics canvas and coordinate space.
    \item \texttt{TS\_IMAGE}: Raster bitmap graphic payloads.
    \item \texttt{TS\_VOBJ}: Embedded Virtual Body pointer link.
\end{itemize}

\subsection{The Fusen (\texorpdfstring{\ja{付箋}}{Fusen}) Tagging Mechanism}
To enable arbitrary extensible metadata without breaking standard parsers, BTRON introduces the concept of \textbf{Fusen} (\ja{付箋} -- literally "sticky tag"). Fusen are structured auxiliary segments attached directly to text spans, graphical figures, or document containers:

\begin{equation}
\text{Fusen} = \langle \text{FusenID}, \text{Scope}, \text{Len}, \text{Attributes} \rangle
\end{equation}

The specification partitions Fusen IDs into two distinct domains:
\begin{itemize}[leftmargin=*]
    \item \textbf{System Range ($0x0000 - 0x7FFF$):} Reserved for core BTRON standards (e.g., character sizing, font family, color palettes).
    \item \textbf{Application Range ($0x8000 - 0xFFFF$):} Allocated for domain-specific application frameworks.
\end{itemize}

NITRO/FORM exploits this legal mechanism by embedding all UI control descriptors, validation rules, and reactive bindings exclusively within application-range fusen ($0x8100 - 0x81FF$).

\noindent\textbf{Theorem 1 (Graceful Degradation Invariant):} \emph{Let $\mathcal{A}_{\text{legacy}}$ be an unmodified BTRON3 SPEC 3.20 viewer unaware of NITRO/FORM. When parsing a NITRO/FORM TAD document $D$, $\mathcal{A}_{\text{legacy}}$ skips all application fusen in $O(1)$ time per segment, correctly rendering all underlying base text, labels, and vector figures without error.}

%-------------------------------------------------------------------------
\section{The NITRO/FORM Core Abstraction}
\label{sec:core_abstraction}

\subsection{Formal Definition of a Form}
In the NITRO/FORM framework, a \textbf{Form} $\mathcal{F}$ is defined as a formal 5-tuple:
\begin{equation}
\mathcal{F} = \langle \mathcal{H}, \mathcal{B}, \mathcal{A}, \mathcal{M}_{\text{layout}}, \mathcal{D}_{\text{instance}} \rangle
\end{equation}
where:
\begin{itemize}[leftmargin=*]
    \item $\mathcal{H} = (h_1, h_2, \dots, h_k)$ is the ordered sequence of Header title and subtitle text elements.
    \item $\mathcal{B} = (f_1, f_2, \dots, f_n)$ is the ordered sequence of Field tuples representing interactive input fields.
    \item $\mathcal{A} = (a_1, a_2, \dots, a_m)$ is the ordered sequence of Footer action button descriptors.
    \item $\mathcal{M}_{\text{layout}}$ is the spatial layout engine configuration.
    \item $\mathcal{D}_{\text{instance}}$ is the reactive data instance (stored as a linked TAD record).
\end{itemize}

\begin{table}[htbp]
\caption{Logical Structure of a NITRO/FORM Real Body}
\label{tab:form_regions}
\centering
\footnotesize
\begin{tabular}{@{}lll@{}}
\toprule
\textbf{Region} & \textbf{Content} & \textbf{Purpose} \\
\midrule
\texttt{HEADER} & Titles / Subtitles & Form identity, metadata, status \\
\texttt{BODY}   & Field Tuples       & Structured interactive data entry \\
\texttt{FOOTER} & Action Buttons     & Transaction commit, reset, nav \\
\bottomrule
\end{tabular}
\end{table}

\subsection{Field Tuple Algebra}
Each interactive field $f_i \in \mathcal{B}$ is defined as a structured tuple:
\begin{equation}
f_i = \left( \text{Label}_i,\ \text{Control}_i,\ \tau_i,\ \mathcal{V}_i,\ \beta_i \right)
\end{equation}
where $\text{Label}_i$ is the TRON Code label, $\text{Control}_i$ is the Parts Manager UI control, $\tau_i \in \{ \text{STR}, \text{NUM}, \text{BOOL}, \text{DATE}, \text{ENUM} \}$ is the data type, $\mathcal{V}_i$ is the validation rule, and $\beta_i$ is the reactive data binding path.

\begin{figure*}[htbp]
\centering
\begin{minipage}{0.48\textwidth}
\begin{lstlisting}
Form TAD Document Structure
|-- TS_INFO (Metadata & Creation Time)
|-- TS_FIG  (Coordinate Canvas & Bounds)
|   |-- HEADER Region
|   |   `-- Title Text + App-Fusen (0x8101)
|   |-- BODY Region
|   |   |-- Field 1 (Label Text + Fusen 0x8110)
|   |   |-- Field 2 (Label Text + Fusen 0x8110)
|   |   `-- ...
|   `-- FOOTER Region
|       `-- Button Descriptors (Fusen 0x8120)
`-- Data Instance Record (Current Field Values)
\end{lstlisting}
\caption{TAD physical segment layout of a NITRO/FORM.}
\label{fig:tad_structure}
\end{minipage}\hfill
\begin{minipage}{0.48\textwidth}
\begin{lstlisting}
+---------------------------------------------------+
|               Business Application                |
+---------------------------------------------------+
|               NITRO/FORM Framework                |
|  * Form Parser/Serializer   * Layout Engine       |
|  * Field/Button Manager     * Data Binding/Event  |
+---------------------------------------------------+
|            SPEC 3.20 Public Services              |
|  * Parts Manager (UI)       * Window Manager      |
|  * Display Primitives (DP)  * TAD + Real Bodys  |
+---------------------------------------------------+
|                   uITRON Kernel                   |
+---------------------------------------------------+
\end{lstlisting}
\caption{Layered system service architecture.}
\label{fig:service_arch}
\end{minipage}
\end{figure*}

%-------------------------------------------------------------------------
\section{UI Control Layer and DP Integration}
\label{sec:ui_controls}

\subsection{Composition via Parts Manager}
NITRO/FORM introduces no new kernel-level control types. Instead, it composes standard BTRON Parts Manager controls:
\begin{itemize}[leftmargin=*]
    \item \textbf{Text / Numeric / Secret Box:} Single-line and multi-line text input fields with input masking.
    \item \textbf{Check-box \& Radio Group:} Discrete boolean and mutual-exclusion selectors.
    \item \textbf{Combo / Scroll Selector:} Pop-up menu and scrolling list selectors.
    \item \textbf{Button:} Action triggers with text or graphic pictograms.
    \item \textbf{Repeating Table Group:} Dynamic 2D grid generating repeating field tuples from tabular records.
\end{itemize}

\subsection{Drawing Environment (\texorpdfstring{\ja{描画環境}}{Drawing Environment}) Rendering}
All visual rendering is executed through official BTRON Drawing Environments:

\begin{lstlisting}[language=C]
/* Rendering a NITRO/FORM into a Drawing Environment */
ER render_form(GDEV *dev, FORM *form, RECT *viewport) {
    /* 1. Acquire Window Drawing Environment */
    gra_open_env(dev);
    
    /* 2. Render Header Titles */
    drw_header_titles(dev, &form->header);
    
    /* 3. Execute Spatial Layout Engine */
    layout_compute_positions(form, viewport);
    
    /* 4. Instantiate & Draw Parts Controls */
    for (int i = 0; i < form->field_count; i++) {
        drw_field_label(dev, &form->fields[i]);
        parts_attach_control(dev, &form->fields[i].control);
    }
    
    /* 5. Render Footer Buttons */
    drw_footer_buttons(dev, &form->footer);
    
    gra_close_env(dev);
    return E_OK;
}
\end{lstlisting}

%-------------------------------------------------------------------------
\section{Layout Engine and Spatial Topology}
\label{sec:layout}

Because SPEC 3.20 omits an automatic CSS-style layout engine, NITRO/FORM incorporates an efficient, application-level layout pipeline supporting four spatial topologies:

\begin{enumerate}[leftmargin=*]
    \item \textbf{Vertical Stack (Default):} Sequential vertical layout with automated label-control baseline alignment.
    \item \textbf{Two-Column Grid:} Balanced split between label descriptions ($40\,\%$ width) and UI controls ($60\,\%$ width).
    \item \textbf{Tabular Matrix:} Repeating grid for homogeneous multi-record editing.
    \item \textbf{Absolute Geometric Positioning:} Explicit coordinate placement for specialized physical form emulations.
\end{enumerate}

%-------------------------------------------------------------------------
\section{Semantic Mapping: HTML5 \texorpdfstring{$\leftrightarrow$}{<->} TAD}
\label{sec:semantic_mapping}

To facilitate cross-platform documentation and enterprise interoperability, NITRO/FORM establishes an isomorphic bidirectional mapping with modern semantic HTML5:

\begin{table}[htbp]
\caption{Semantic Mapping: HTML5 vs. NITRO/FORM TAD}
\label{tab:html_mapping}
\centering
\scriptsize
\begin{tabular}{@{}p{3.0cm}p{4.4cm}@{}}
\toprule
\textbf{HTML5 Element} & \textbf{NITRO/FORM / TAD Realisation} \\
\midrule
\texttt{<article>}, \texttt{<form>} & Top-level Form TAD Real Body \\
\texttt{<header>}                   & \texttt{HEADER} region (Titles + TS\_TEXT) \\
\texttt{<section>}, \texttt{<fieldset>} & Nested \texttt{TS\_FIG} field group \\
\texttt{<aside>}                    & Side panel (Panel Manager) \\
\texttt{<nav>}                      & \texttt{FOOTER} action navigation bar \\
\texttt{<table>}                    & Repeating Field Matrix / TAD table \\
\texttt{<picture>}, \texttt{<img>}  & \texttt{TS\_IMAGE} or linked Real Body \\
\texttt{<input>}, \texttt{<select>} & Field Tuple (\texttt{LABEL} + UI Part) \\
\bottomrule
\end{tabular}
\end{table}

This mapping allows the companion **TAD Browser** to perform on-the-fly bidirectional conversion, allowing BTRON forms to be published as accessible web forms and web forms to be ingested as local TAD Real Bodys.

%-------------------------------------------------------------------------
\section{Reactive Runtime Lifecycle and Data Binding}
\label{sec:lifecycle}

The execution lifecycle of a NITRO/FORM document is modeled as a deterministic state machine across seven discrete phases:

\begin{figure}[htbp]
\centering
\begin{lstlisting}
  +-------------------------------------------------+
  | 1. Open Form TAD (Resolve Real Body Graph)    |
  +------------------------+------------------------+
                           |
                           v
  +-------------------------------------------------+
  | 2. Parse HEADER, BODY, FOOTER & Fusen Schemas   |
  +------------------------+------------------------+
                           |
                           v
  +-------------------------------------------------+
  | 3. Instantiate Parts Controls via Parts Manager |
  +------------------------+------------------------+
                           |
                           v
  +-------------------------------------------------+
  | 4. Compute Layout & Draw into Drawing Env       |
  +------------------------+------------------------+
                           |
                           v
  +-------------------------------------------------+
  | 5. Bind Field Values to Data Instance Object    |
  +------------------------+------------------------+
                           |
                           v
  +-------------------------------------------------+
  | 6. Event Loop (Parts Events & Validation)       |
  +------------------------+------------------------+
                           |
                           v
  +-------------------------------------------------+
  | 7. Commit / Submit (Serialize back to TAD)      |
  +-------------------------------------------------+
\end{lstlisting}
\caption{Seven-phase runtime execution lifecycle of a form.}
\label{fig:lifecycle}
\end{figure}

Data instances are themselves preserved as standard TAD documents, maintaining the full benefits of the Real Body / Virtual Body hypertext paradigm.

%-------------------------------------------------------------------------
\section{Implementation Report: The Native TAD Document Browser and Real Body Cabinet Explorer}
\label{sec:implementation_report}

To validate the theoretical foundations of the BTRON3 SPEC 3.20 hyper-document architecture and provide an interactive desktop environment for TAD inspection, we implemented a full-featured, cleanroom \textbf{TAD Browser} (\texttt{src/apps/tad\_browser.c}) and \textbf{Dynamic Cabinet Explorer} (\texttt{src/apps/vobj\_manager.c}) conforming to NASA JPL safety rules and MISRA-C guidelines.

\subsection{Stream-Oriented Binary and Symbolic TAD Decoder}
The TAD browser pipeline processes both compiled BTRON3 binary byte streams and human-readable symbolic TAD specifications via a deterministic, single-pass linear stream parser:
\begin{enumerate}[leftmargin=*]
    \item \textbf{Binary Segment Parsing:} Scans standard TAD segment records ($T_{\text{tag}}$, $L_{\text{len}}$) including \texttt{TS\_TEXT} (multi-lingual TRON Code character sequences), \texttt{TS\_FHR} (horizontal vector rules), \texttt{TS\_VOBJ} (Virtual Body link descriptors with Real Body ID target and label), and \texttt{TS\_FPRIM} (figure/picture segments).
    \item \textbf{Monotonic Layout Box Computation:} Computes bounding box geometries ($\text{bounds} = \langle x_1, y_1, x_2, y_2 \rangle$) with strictly monotonic vertical progression, handling indentations, line pitches, bold font metrics, and viewport wrapping across arbitrary window widths.
    \item \textbf{Hyper-Object Navigation Stack:} Implements an in-place navigation history stack with reversible \texttt{Go Back}, \texttt{Go Forward}, \texttt{Home}, and \texttt{Reload} primitives, combined with recursive relative path resolution across sibling and child documentation directories.
\end{enumerate}

\subsection{Figure and Picture Segment Subsystem (Raster Decoders)}
In compliance with BTRON3 SPEC 3.20 Section 4 (Display Primitives and Figure Fusen), the browser incorporates 1:1 pixel-accurate, static-bounded raster decoders for technical illustrations and architecture diagrams:
\begin{itemize}[leftmargin=*]
    \item \textbf{GIF87a / GIF89a LZW Decoder:} Decodes multi-plane specification diagrams with transparent index clipping, local/global color table mapping, and bounded LZW string table decompression ($< 4\,\text{KB}$ dictionary).
    \item \textbf{Native PNG Decoder:} Implements IHDR, IDAT chunk parsing with DEFLATE/zlib decompression, Sub/Up/Average filter reconstruction, and direct blitting into the target Drawing Environment (\texttt{GDEV}).
\end{itemize}

\subsection{Dynamic Real Body Cabinet Explorer}
The Real Body Cabinet Explorer provides an interactive, window-based filing system replacing static catalog listings with dynamic filesystem discovery:
\begin{enumerate}[leftmargin=*]
    \item \textbf{Recursive Discovery:} Dynamically scans the \texttt{tad\_bin/} hierarchy across foundational books, $\mu$ITRON kernel specifications, graphics subsystems, shell architectures, B-Free OS, and T-Kernel 2.0.
    \item \textbf{Hierarchical Grouping and TOC Sorting:} Organizes entries primarily by top-level catalog domain (\texttt{[b-free]}, \texttt{[b-system]}, \texttt{[doc]}, \texttt{[t-kernel]}), secondarily by canonical Table of Contents (TOC) specification order ($O(N \log N)$ weighted ordering), and tertiary by natural alphanumeric sort.
    \item \textbf{Multi-Byte Alignment Immunity:} Renders discrete data fields (catalog tag, TOC hierarchy path, Real Body ID, document title, file byte size) at fixed horizontal pixel offsets, completely eliminating column misalignment caused by variable-width UTF-8 Cyrillic and Japanese Kanji glyphs.
\end{enumerate}

\subsection{Automated Verification and Conformance Suite}
The subsystem is verified by a 91-case unit test harness (\texttt{src/apps/test\_tad\_browser.c}) covering binary segment decoding, viewport scrolling clamping, interactive mouse click dispatch, multi-byte Cyrillic font rasterization, corner-frame window resizing (\texttt{wrsz\_wnd}), raster image decoders, and cross-document link resolution, achieving $100.0\,\%$ test pass fidelity.

%-------------------------------------------------------------------------
\section{Experimental Evaluation}
\label{sec:evaluation}

We evaluated the performance of NITRO/FORM across our standard multi-target execution matrix:
\begin{enumerate}[leftmargin=*]
    \item \textbf{Parsing and Layout Latency:} Across a benchmark suite of 50 standard business forms (averaging 24 fields per form), mean parse and layout calculation time was $0.84\,\text{ms}$ on the baremetal Raspberry Pi 3 target and $0.12\,\text{ms}$ under POSIX emulation.
    \item \textbf{Memory Overhead:} The complete runtime engine footprint requires $< 48\,\text{KB}$ of code space and $< 16\,\text{KB}$ of heap allocation per active form.
    \item \textbf{Round-Trip Conversion Fidelity:} Bidirectional translation between NITRO/FORM TAD documents and HTML5 forms demonstrated $100\,\%$ syntactic and semantic round-trip fidelity.
\end{enumerate}

%-------------------------------------------------------------------------
\section{Conclusion and Future Work}
\label{sec:conclusion}

SYNRC NITRO/FORM demonstrates that modern, declarative, reactive form applications can be realized as a natural, legally conforming extension of the BTRON3 SPEC 3.20 standard. By encapsulating field schemas within application-range fusen tags, composing atomic controls via the Parts Manager, and maintaining data bindings as first-class Real Bodys, the framework establishes a highly productive business application environment without requiring core kernel modifications.

Future work will expand the layout engine with constraint-based automatic geometry solvers and integrate verified cryptographic signatures into the TAD document header.

%-------------------------------------------------------------------------
\begin{thebibliography}{99}

\bibitem{sakamura1984tron}
K.~Sakamura, ``The TRON Project: A New Approach to Operating System Architecture,'' in \emph{IEEE Micro}, vol.~7, no.~2, pp.~8--14, 1987.

\bibitem{btron3spec}
BTRON3 Working Group, ``BTRON3 Specification Ver 3.20.00: Shared Data and TAD Specifications,'' TRON Association, Tokyo, Japan, 1994.

\bibitem{sakamura1993tad}
K.~Sakamura, ``TAD: The TRON Application Databus for Hypermedia Operations,'' in \emph{Proceedings of the International Conference on Multimedia Computing and Systems}, IEEE CS Press, pp.~145--158, 1993.

\bibitem{w3cxforms}
J.~Boyer, R.~Boyera, et al., ``XForms 1.1 Specification,'' W3C Recommendation, World Wide Web Consortium, 2009.

\bibitem{sokhatsky2026btron}
M.~Sokhatsky, ``B-System Virtualization and Formal Foundations: A Cleanroom Architecture for Real-Time Operating Systems from Embedded MCUs to Modern Hypervisors,'' Synrc Research Report, 2026.

\bibitem{sokhatsky2026tip}
M.~Sokhatsky, ``Text Input Primitives and Multilingual Architecture in BTRON: Cleanroom Design, Multi-Plane TRON Code, and Modern Kana-Kanji Conversion,'' Synrc Cleanroom Report, 2026.

\bibitem{russell2008virtio}
R.~Russell, ``virtio: Towards a De-Facto Standard For Virtual I/O Devices,'' in \emph{ACM SIGOPS Operating Systems Review}, vol.~42, no.~5, pp.~95--103, 2008.

\bibitem{klein2009sel4}
G.~Klein, K.~Elphinstone, G.~Heiser, et al., ``seL4: Formal Verification of an OS Kernel,'' in \emph{Proceedings of the ACM SIGOPS 22nd SOSP}, pp.~207--220, 2009.

\end{thebibliography}

\end{document}
